%TODO: cite the correct code
\label{sec:impl}
NDNSF was implemented using \texttt{ndn-cxx}~\cite{ndn-cxx}, NDN-SVS~\cite{ndnsvs-github}, and NAC-ABE~\cite{dulal2023nac-abe}.

NDNSF validates certificate chains against fixed trust anchors, checks certificate validity periods, and enforces configured application trust rules. Cached certificate validation expires with the validated chain.

For NAC-ABE, cached keys avoid repeated retrieval and decryption. Small publications can also be carried directly in Sync Interest Application Parameters, avoiding a separate publication fetch. Regular-expression subscriptions filter message types and Request service names; receivers then check the recipient, service and invocation bindings.

Peers can reference immutable, producer-signed security profiles by exact Data name under the Producers' prefixes. Receivers fetch and validate uncached profiles, reuse their encryption certificates and capabilities within certificate validity, and retain authorization and request-binding checks for each invocation.

SVS byte-publication preparation is synchronous by default; optional bounded workers can prepare publications asynchronously. Network operations remain on the Face event loop. The evaluation measures the combined invocation path rather than each optimization separately.

Request-scoped protection is enabled by default for Controller-managed calls. With custom capability settings, input and result isolation depends on retaining request-scoped protection.

A two-provider MiniNDN functional experiment verified protected nonempty input delivery and results for both selected providers; in a separate call, a nonselected provider did not obtain decrypted input or usable Selection keys. The performance experiments below select one provider per invocation.

%NDNSF was implemented using ndn-cxx~\cite{ndn-cxx}, NDN-SVS~\cite{ndnsvs-github}, and NAC-ABE~\cite{dulal2023nac-abe}. During development, we observed that the baseline NAC-ABE and NDN-SVS implementations did not adequately meet performance expectations, and we therefore introduced several optimizations.

%For NAC-ABE, we implemented a key-caching mechanism to eliminate repeated key retrieval and decryption, reducing decryption latency by approximately one RTT per operation. For NDN-SVS, messages smaller than 512~bytes are embedded in the Sync Interest’s Application Parameters, reducing the publication-to-reception delay by one RTT. In addition, to ensure that users and providers receive only relevant messages, we introduced a regular-expression-based subscription mode into NDN-SVS.
%For example, the user \path{/UAVNET/Drone3} requesting the \path{ImageProc/YOLO12} service uses the following two regular expressions in its subscriptions: \path{"<>+/<NDNSF>/<ACK>/<UAVNET>/<Drone3>/<ImageProc>/<YOLO12>/<>"} for ACK messages and \path{"<>+/<NDNSF>/<RESPONSE>/<UAVNET>/<Drone3>/<ImageProc>/<YOLO12>/<>"} for Response messages. On the other hand, the provider \path{/UAVNET/Drone1}, which offers the service \path{/UAVNET/ImageProc/Yolov12}, subscribes to Request messages using \path{"<>+<NDNSF><REQUEST><UAVNET><ImageProc><Yolov12><>"} and to Selection messages using \path{"<>+<NDNSF><SELECTION><UAVNET><Drone1><ImageProc><Yolov12><>"}.

%Finally, under high-frequency NDNSF workloads, the original blocking APIs and heavily single-threaded design of NDN-SVS became a major bottleneck, since most synchronization and processing logic executed on the main Face event loop. To address this issue, we redesigned Sync Interest handling and publication into an asynchronous pipeline. CPU-intensive tasks such as state-vector processing, decoding, comparison, and metadata preparation are offloaded to a bounded worker pool, while all NDN Face and \texttt{io\_context} operations remain on the main NDN thread to preserve thread safety. In addition, we introduced a nonblocking asynchronous publish path that prevents application-level request handling from stalling the Face event loop. These optimizations significantly improve NDNSF scalability and responsiveness under high-frequency request workloads.

%To demonstrate how NDNSF can be used to build distributed applications and evaluate its performance, we developed a UAV system demo in which the ground station invokes the Flight Control service provided by the drone. This demo utilizes NDNSF for inter-node communication and access control, and once the drone's onboard computer verifies the command, it uses the MAVSDK API to control the UAV flight controller, thereby enabling autonomous operation of the UAV.
